\chapter{Charts and sparklines}
\label{ch:charts}

\begin{aside}
A 80-cell terminal can show a graph. Not as detailed as a
1920-pixel display, but enough to spot a trend. Sparklines,
bar charts, line graphs --- all are widgets in \maya{}.
\end{aside}

\section{Sparkline}

A compact line graph in one row using block characters:

\begin{lstlisting}
constexpr const char* blocks[] = {
    "\u2581", "\u2582", "\u2583", "\u2584",
    "\u2585", "\u2586", "\u2587", "\u2588"
};

void sparkline(Painter& p, span<float> data, Rect r) {
    auto [lo, hi] = std::minmax_element(...);
    for (size_t i = 0; i < data.size(); ++i) {
        int level = scale(data[i], *lo, *hi, 0, 7);
        p.put(r.x + i, r.y, Cell{blocks[level]});
    }
}
\end{lstlisting}

8 levels of fill per cell give decent vertical resolution
in one row.

\section{Bar chart}

A multi-row bar chart uses the same block characters,
stacking vertically. Each column is a value.

\section{Line chart}

For smooth lines across multiple rows, use braille
characters. Each braille cell is a 2x4 dot grid; a line
crosses cells by lighting the appropriate dots.

\maya{}'s \code{line\_chart} widget does this.

\section{Axes and labels}

A chart with axes needs labelled gridlines. Render the
labels at fixed positions; ensure they don't overlap.

\section{Colours}

Multiple series get different colours. Avoid red-green
colour blindness traps; use distinct hues.

\section{Real-time updates}

A live chart appends a new data point per frame. The
chart scrolls left, dropping the oldest.

\section{Logarithmic}

For data with large dynamic range, log scale. Provide a
toggle.

\section{Recap}

\begin{recap}
\begin{itemize}[leftmargin=*]
  \item Sparkline = one row, 8 fill levels per cell.
  \item Bar chart = many rows of fills.
  \item Line chart = braille for sub-cell precision.
  \item Real-time scroll for live data.
\end{itemize}
\end{recap}

\section*{Workshop}

\begin{enumerate}[leftmargin=*]
  \item Build a sparkline showing CPU usage.
  \item Add a real-time line chart of network throughput.
  \item Stack two series with different colours.
\end{enumerate}
